\section{Methodology}
\label{sec:method}

The strategy is re-estimated at the close of every trading day $t$ using only returns up to and
including day $t$. \Cref{tab:params} lists the parameters and their baseline values, which
follow AL. \Cref{alg:daily} summarises one day of the procedure.

\begin{table}[ht]
\centering
\caption{Strategy parameters. Baseline values follow \citet{avellaneda2010}.}
\label{tab:params}
\begin{tabular}{llr}
\toprule
Symbol & Meaning & Baseline \\
\midrule
$M$ & correlation estimation window (days) & 252 \\
$m$ & number of eigenportfolios (factors) & 15 \\
$L$ & residual estimation window (days) & 60 \\
$\kappa_{\min}$ & minimum mean-reversion speed (per year) & $252/30 \approx 8.4$ \\
$\bar s_{\text{open}}$ & entry threshold & 1.25 \\
$\bar s_{\text{cl}}$ & exit threshold, longs & 0.50 \\
$\bar s_{\text{cs}}$ & exit threshold, shorts & 0.75 \\
$q$ & notional per position (fraction of capital) & 0.01 \\
$c$ & one-way transaction cost & 5~bps \\
$\delta$ & no-trade band (fraction of capital) & 0 (off) \\
\bottomrule
\end{tabular}
\end{table}

\subsection{Eigenportfolios}
\label{sec:method:pca}

Let $\mathcal U_t$ be the set of $N_t$ stocks with a complete return history over the window
$\{t-M+1,\dots,t\}$. For each such stock compute the sample mean $\bar R_i$ and standard
deviation $\sigma_i$ over the window, and the standardised returns
$Y_{is} = (R_{is} - \bar R_i)/\sigma_i$. The empirical correlation matrix is
\begin{equation}
  \rho_{ij} = \frac{1}{M-1} \sum_{s=t-M+1}^{t} Y_{is} Y_{js} .
\end{equation}
Let $\lambda_1 \ge \lambda_2 \ge \dots \ge \lambda_{N_t}$ be its eigenvalues and
$\bm v^{(1)},\dots,\bm v^{(N_t)}$ the corresponding unit eigenvectors. The first $m$
eigenvectors define the \emph{eigenportfolios}, with dollar weights
\begin{equation}
  Q_i^{(j)} = \frac{v_i^{(j)}}{\sigma_i}, \qquad j = 1,\dots,m,
  \label{eq:eigenportfolio}
\end{equation}
and returns $F_{js} = \sum_i Q_i^{(j)} R_{is}$. Dividing by $\sigma_i$ maps the eigenvector,
which lives in the space of standardised returns, back to dollar holdings, so volatile stocks
receive smaller weights. Eigenvectors are only defined up to sign. We choose the sign that
makes $\sum_i v_i^{(j)} > 0$. This does not affect the strategy but makes the first factor read
as a long-only ``market'' portfolio.

\subsection{Factor regression}
\label{sec:method:regression}

For each stock $i \in \mathcal U_t$, the last $L$ daily returns are regressed on the
contemporaneous eigenportfolio returns:
\begin{equation}
  R_{is} = \alpha_i + \sum_{j=1}^{m} \beta_{ij} F_{js} + \varepsilon_{is},
  \qquad s = t-L+1,\dots,t .
  \label{eq:regression}
\end{equation}
The correlation matrix uses a long window ($M=252$) because the factors need many observations
to be estimated stably. The regression uses a short window ($L=60$) because factor loadings
drift over time. All $N_t$ regressions share the same design matrix, so they are solved with a
single least-squares call (\cref{sec:impl}).

\subsection{Residual process}
\label{sec:method:ou}

The cumulative residual
\begin{equation}
  X_{ik} = \sum_{s=t-L+1}^{t-L+k} \varepsilon_{is}, \qquad k = 1,\dots,L,
\end{equation}
measures how far stock $i$ has drifted from the path implied by its factor exposures. Following
AL, it is modelled as an Ornstein--Uhlenbeck process,
\begin{equation}
  dX_i = \kappa_i (m_i - X_i)\, dt + \sigma_i^{\text{OU}}\, dW_i ,
  \label{eq:ou}
\end{equation}
which reverts to the level $m_i$ at speed $\kappa_i$. Sampled daily, \cref{eq:ou} is exactly an
AR(1) process (\cref{app:ou}),
\begin{equation}
  X_{i,k+1} = a_i + b_i X_{ik} + \zeta_{i,k+1},
  \label{eq:ar1}
\end{equation}
which is fitted by least squares. The OU parameters follow as
\begin{equation}
  \kappa_i = -252 \ln b_i, \qquad
  m_i = \frac{a_i}{1-b_i}, \qquad
  \sigma_{\text{eq},i} = \sqrt{\frac{\Var(\zeta_i)}{1-b_i^2}},
  \label{eq:ou-params}
\end{equation}
where $\sigma_{\text{eq},i}$ is the standard deviation of the stationary distribution of
$X_i$. If $b_i \notin (0,1)$, the fitted process is not mean-reverting and the stock is not
eligible on that day. The stock is also ineligible if its estimated mean-reversion time
$1/\kappa_i$ exceeds $1/\kappa_{\min}$, about 30 trading days, i.e.\ half the estimation window.
Reversion slower than that cannot be estimated reliably from $L=60$ observations.

\subsection{The s-score}
\label{sec:method:sscore}

For eligible stocks, the \emph{s-score} is the distance of the current cumulative residual from
its equilibrium, measured in equilibrium standard deviations:
\begin{equation}
  s_i = \frac{X_{iL} - \tilde m_i}{\sigma_{\text{eq},i}}, \qquad
  \tilde m_i = m_i - \frac{1}{|\mathcal E_t|}\sum_{k \in \mathcal E_t} m_k ,
  \label{eq:sscore}
\end{equation}
where $\mathcal E_t$ is the set of eligible stocks. Centring the equilibrium levels across
stocks removes a common bias in the estimated $m_i$ that AL attribute to the regression
intercepts. Because \cref{eq:regression} includes an intercept, the residuals sum to zero over
the window, so $X_{iL} = 0$ exactly and $s_i = -\tilde m_i/\sigma_{\text{eq},i}$. A strongly
negative s-score means the stock is cheap relative to its factors and is expected to rise
relative to them. A strongly positive s-score means the reverse.

\subsection{Trading rule}
\label{sec:method:rule}

Each stock is in one of three states: long ($+1$), flat ($0$) or short ($-1$). States are
updated daily by the rule in \cref{tab:rule}. Exits are applied before entries, so a position can
reverse on a single day. The gap between the entry and exit thresholds (hysteresis) prevents
positions from being repeatedly opened and closed as the s-score fluctuates around a single
threshold.

\begin{table}[ht]
\centering
\caption{Trading rule on the s-score.}
\label{tab:rule}
\begin{tabular}{lll}
\toprule
Current state & Condition & New state \\
\midrule
flat & $s_i < -\bar s_{\text{open}}$ & long \\
flat & $s_i > +\bar s_{\text{open}}$ & short \\
long & $s_i > -\bar s_{\text{cl}}$ & flat \\
short & $s_i < +\bar s_{\text{cs}}$ & flat \\
any & stock ineligible on day $t$ & flat \\
\bottomrule
\end{tabular}
\end{table}

\subsection{Portfolio construction and hedging}
\label{sec:method:portfolio}

Let $\theta_{it} \in \{-1,0,1\}$ be the state of stock $i$ and $w_i = q\,\theta_{it}$ the
dollar position in the stock. To isolate the residual, each position is hedged by selling
$\beta_{ij} w_i$ dollars' worth of eigenportfolio $j$ for every $j$. One unit of eigenportfolio
$j$ holds $Q_k^{(j)}$ dollars of stock $k$, so the hedges can be netted back into single-stock
holdings. Writing $Q \in \mathbb R^{N\times m}$ and $B \in \mathbb R^{N\times m}$ for the
matrices of eigenportfolio weights and betas, the net holdings are
\begin{equation}
  \bm w^{\text{net}} = \bm w - Q B^\top \bm w .
  \label{eq:netting}
\end{equation}
Within the estimation window the portfolio return is
$\sum_k w^{\text{net}}_k R_k = \sum_i w_i (R_i - \sum_j \beta_{ij} F_j) = \sum_i w_i(\alpha_i + \varepsilon_i)$:
the factor exposures cancel by construction and only the residuals remain.

\paragraph{Daily versus frozen hedges.}
In AL's specification, \cref{eq:netting} is re-evaluated every day with that day's $Q$ and $B$,
so the hedge of every open position is continuously updated. We also implement a
\emph{frozen} hedge. The stock-space hedge of position $i$,
$\bm h_i = -w_i\, Q^{(\tau_i)} \bm\beta_i^{(\tau_i)}$, is computed once on the entry day
$\tau_i$ and held unchanged until the position is closed. Net holdings are then
$\bm w + \sum_i \bm h_i$. The trading signals are identical in both modes. Only the hedge
differs.

\subsection{Cost-reduction mechanisms}
\label{sec:method:costs}

\paragraph{Volume-adjusted returns.}
AL's ``trading time'' variant replaces $R_{is}$ in \cref{eq:regression} by
$R_{is}\,\langle V_i\rangle_s / V_{is}$, where $V_{is}$ is the share volume and
$\langle V_i\rangle_s$ its trailing 10-day mean ending on day $s$. Returns on heavy-volume days,
which are more likely to reflect news, are shrunk, and returns on light-volume days are
amplified. Only the signal uses adjusted returns. The factors and the P\&L use actual returns.

\paragraph{No-trade band.}
Given the target holdings $w^\ast_{kt}$ from \cref{eq:netting} and yesterday's holdings
$w_{k,t-1}$, each stock is traded only if its required adjustment exceeds the band $\delta$:
\begin{equation}
  w_{kt} =
  \begin{cases}
    w^\ast_{kt} & \text{if } |w^\ast_{kt} - w_{k,t-1}| > \delta, \\
    w_{k,t-1} & \text{otherwise.}
  \end{cases}
  \label{eq:band}
\end{equation}
This is a simple form of the no-trade region that is optimal under proportional costs
\citep{davis1990}. It removes the many small hedge adjustments spread across all stocks, at the
price of a small tracking error to the target.

\paragraph{Fewer factors.}
Every eigenportfolio holds a position in every stock. So each factor hedged adds a basket trade
across the entire universe whenever a position is opened or closed. Hedging fewer factors
reduces this cost but leaves more systematic risk in the residual. This trade-off is studied
in \cref{sec:results:factors}.

\subsection{Backtest accounting}
\label{sec:method:accounting}

Holdings determined at the close of day $t$ earn the returns of day $t+1$, so no information from
day $t+1$ is used to form them. Capital is fixed at one unit, without compounding, as in AL.
Daily P\&L, as a fraction of capital, is
\begin{align}
  \text{gross}_{t+1} &= \sum_k w_{kt}\, R_{k,t+1}, \\
  \text{cost}_{t+1} &= c \sum_k \left| w_{kt} - w_{k,t-1} \right|, \\
  \text{net}_{t+1} &= \text{gross}_{t+1} - \text{cost}_{t+1}.
\end{align}
The quantity $\text{TO}_t = \sum_k |w_{kt} - w_{k,t-1}|$ is the daily one-way turnover. A missing
next-day return of a held stock (for example after a delisting) is treated as zero.

Annualised return and volatility are $252$ times the mean and $\sqrt{252}$ times the standard
deviation of daily net P\&L, and the Sharpe ratio is their ratio. Because the strategy is
market-neutral and self-financing apart from margin, no risk-free rate is subtracted. We also
report the \emph{break-even cost}
\begin{equation}
  c^\ast = \frac{\overline{\text{gross}}}{\overline{\text{TO}}},
  \label{eq:breakeven}
\end{equation}
the one-way cost at which net P\&L would be zero. It measures the gross edge per unit of capital
traded and is the most direct summary of whether a variant can survive costs.

\begin{algorithm}[ht]
\caption{One trading day $t$ of the strategy}
\label{alg:daily}
\begin{algorithmic}[1]
\State $\mathcal U_t \gets$ stocks with complete returns over $[t-M+1, t]$
\State $Q \gets$ eigenportfolio weights from the top $m$ eigenvectors of $\rho$ \Comment{\cref{eq:eigenportfolio}}
\State $F \gets R_{[t-L+1,t],\,\mathcal U_t}\, Q$ \Comment{factor returns, $L\times m$}
\State $(\bm\alpha, B, \bm\varepsilon) \gets$ least squares of $R$ on $[\bm 1, F]$ \Comment{\cref{eq:regression}}
\State $X \gets \operatorname{cumsum}(\bm\varepsilon)$; fit AR(1); compute $\kappa, m, \sigma_{\text{eq}}$ \Comment{\cref{eq:ar1,eq:ou-params}}
\State $\bm s \gets$ s-scores, NaN where $\kappa < \kappa_{\min}$ \Comment{\cref{eq:sscore}}
\State $\bm\theta_t \gets$ update states with the rule of \cref{tab:rule}
\State $\bm w^\ast \gets q\bm\theta_t - Q B^\top q\bm\theta_t$ \Comment{or frozen hedges}
\State $\bm w_t \gets$ apply the no-trade band to $\bm w^\ast$ \Comment{\cref{eq:band}}
\State earn $\bm w_t^\top R_{t+1}$; pay $c\,\|\bm w_t - \bm w_{t-1}\|_1$
\end{algorithmic}
\end{algorithm}
